\documentclass{exam}
\usepackage{ModeloIFBA}
%\usepackage[alf]{abntex2cite}
\usepackage{multicol}
\usepackage{adjustbox}
\usepackage{multicol}

\setlength\columnsep{30pt}


\renewcommand{\curso}{Tecnologia em Análise e Desenvolvimento de Sistemas }
\renewcommand{\materia}{Programação Orientada a Objetos}
\renewcommand{\turma}{2024.1}
\renewcommand{\professor}{Leandro Costa Souza}
\renewcommand{\avaliacao}{Prova Final}
\renewcommand{\dataavaliacao}{08/10/2024}


\newcommand{\duracao}{1h 40min}
\newcounter{solcount}
\setcounter{solcount}{1}
\newcommand\CC{%
    \CorrectChoice\label{sol:\arabic{solcount}}
    \stepcounter{solcount}
}

\newcommand\printmyanswers{%
\section*{\centering Answers Sheet}
    \foreach \x in {1,...,\thequestion}{%
    \noindent
        \x.~\ref{sol:\x}\par
    }
}

\begin{document}
\begin{adjustbox}{max width=\textwidth}
\begin{tabularx}{\textwidth}{rrXrl}
\multirow{5}{*}{\ifbalogo{1.7in}} & \textbf{Curso:} & \multicolumn{3}{l}{ \curso } \\ 
& \textbf{Matéria:} & \materia \space & \textbf{Turma:} & \small{\turma}  \\ 
& \textbf{Professor:} & \professor & & \\ 
& \textbf{Aluno:} & \hrulefill & \textbf{Data:} & \dataavaliacao \\
\end{tabularx}
\end{adjustbox}
\begin{tabularx}{\textwidth}{XcX}
 & \textbf{\avaliacao} &  \\ \hline
\end{tabularx}	
	\input{orientacoes}


	%\begin{center}
	%	\gradetable[v][questions] 
	%\end{center}
	\begin{questions}
\begin{multicols}{2}
    \question[1] Qual método da classe Thread é usado para iniciar a execução de uma nova thread?
\begin{choices}
 \CC start()
 \choice run()
 \choice execute()
 \choice begin()
\end{choices}

\question[1] Qual é a principal diferença entre a interface Runnable e a classe Thread?
\begin{choices}
 \choice Runnable é uma classe abstrata, enquanto Thread é uma interface.
 \choice Runnable permite herança múltipla, enquanto Thread não.
 \CC Runnable define um único método run(), enquanto Thread é uma classe completa que pode ser estendida.
 \choice Thread define um único método run(), enquanto Runnable é uma classe completa que pode ser estendida.
\end{choices}

\question[1] Qual das seguintes opções é verdadeira sobre o uso de synchronized em métodos?
\begin{choices}
 \choice Ele permite que múltiplas threads acessem o método simultaneamente.
 \choice Ele impede que existam múltiplas instancias do objeto para garantir uma única thread.
 \choice Ele impede que qualquer thread acesse o método.
 \CC Ele permite que apenas uma thread acesse o método por vez, mas não impede outras threads de acessar outros métodos da mesma classe.
\end{choices}

\question[1] Qual das seguintes opções é verdadeira sobre exceções verificadas em Java?
\begin{choices}
 \choice Elas são verificadas em tempo de execução.
 \choice Elas não precisam ser declaradas ou tratadas.
 \CC Elas são verificadas em tempo de compilação.
 \choice Elas são subclasses de RuntimeException.
\end{choices}

\question[1] Qual das seguintes opções é verdadeira sobre exceções não verificadas em Java?
\begin{choices}
 \CC Elas são verificadas em tempo de execução.
 \choice Elas são verificadas em tempo de compilação.
 \choice Elas precisam ser declaradas ou tratadas.
 \choice Elas são subclasses de Exception.
\end{choices}

\question[1] Qual palavra-chave é usada para definir que uma exceção personalizada pode ser lançada por um método em Java?
\begin{choices}
 \choice throws
 \CC throw
 \choice exception
 \choice custom
\end{choices}

\question[1] Qual das seguintes opções é verdadeira sobre o bloco finally em Java?
\begin{choices}
 \choice Ele é executado apenas se uma exceção for lançada.
 \choice Ele é executado apenas se nenhuma exceção for lançada.
 \CC Ele é executado sempre, independentemente de uma exceção ser lançada ou não.
 \choice Ele não pode ser usado com blocos try-catch.
\end{choices}

\question[1] Qual das seguintes opções é verdadeira sobre a palavra-chave throw em Java?
\begin{choices}
 \choice Ela é usada para declarar uma exceção.
 \CC Ela é usada para lançar uma exceção.
 \choice Ela é usada para capturar uma exceção.
 \choice Ela é usada para finalizar um bloco try-catch.
\end{choices}

\question[1] Qual das seguintes opções é verdadeira sobre a classe Throwable em Java?
\begin{choices}
 \choice Ela é a superclasse de todas as exceções verificadas.
 \choice Ela é a superclasse de todas as exceções não verificadas.
 \CC Ela é a superclasse de todas as exceções e erros.
 \choice Ela não pode ser estendida por outras classes.
\end{choices}

\question[1] Qual das seguintes opções é verdadeira sobre a classe Exception em Java?
\begin{choices}
 \CC Ela é a superclasse de todas as exceções verificadas.
 \choice Ela é a superclasse de todas as exceções não verificadas.
 \choice Ela é a superclasse de todos os erros.
 \choice Ela não pode ser estendida por outras classes.
\end{choices}

\question[1] Qual das seguintes opções é verdadeira sobre a comunicação entre camadas em uma arquitetura de três camadas?
\begin{choices}
 \choice A camada de apresentação se comunica diretamente com a camada de dados.
 \choice A camada de serviço se comunica diretamente com a camada de apresentação.
 \choice A camada de dados se comunica diretamente com a camada de apresentação.
 \CC A camada de apresentação se comunica com a camada de serviço, que por sua vez se comunica com a camada de dados.
\end{choices}

\question[1] Qual das seguintes opções é uma vantagem da separação em camadas?
\begin{choices}
 \choice Aumenta a complexidade do código.
 \choice Dificulta a manutenção do código.
 \CC Facilita a manutenção e a escalabilidade do código.
 \choice Reduz a modularidade do código.
\end{choices}

\question[1] Qual das seguintes opções é um exemplo de tecnologia comumente usada na camada de apresentação?
\begin{choices}
 \choice JDBC
 \choice JPA
 \choice Hibernate
 \CC JavaFX
\end{choices}

\question[1] Qual palavra-chave é usada para herdar uma classe em Java?
\begin{choices}
 \choice inherits
 \CC extends
 \choice implements
 \choice inheritsFrom
\end{choices}

\question[1] Qual das seguintes opções é verdadeira sobre herança múltipla em Java?
\begin{choices}
 \choice Java suporta herança múltipla de classes.
 \CC Java suporta herança múltipla de interfaces.
 \choice Java não suporta herança múltipla de classes ou interfaces.
 \choice Java suporta herança múltipla de classes e interfaces.
\end{choices}

\question[1] Qual das seguintes opções é verdadeira sobre o uso da palavra-chave super em Java?
\begin{choices}
 \choice Ela é usada para chamar o construtor da classe filha.
 \choice Ela é usada para chamar métodos estáticos da classe pai.
 \choice Ela é usada para chamar métodos privados da classe pai.
 \CC Ela é usada para chamar o comportamento definido da classe pai.
\end{choices}

\question[1] Qual das seguintes opções é verdadeira sobre a classe Object em Java?
\begin{choices}
 \CC Ela é a superclasse de todas as classes em Java.
 \choice Ela é a subclasse de todas as classes em Java.
 \choice Ela não pode ser estendida por outras classes.
 \choice Ela não pode ser instanciada.
\end{choices}

\question[1] Qual das seguintes opções é verdadeira sobre o uso de herança em Java?
\begin{choices}
 \choice Uma classe pode herdar de múltiplas classes.
 \CC Uma classe pode herdar de apenas uma classe.
 \choice Uma classe não pode herdar de outra classe.
 \choice Uma classe pode herdar de múltiplas interfaces e classes.
\end{choices}

\question[1] Qual das seguintes opções é verdadeira sobre o uso de herança para reutilização de código?
\begin{choices}
 \choice Herança impede a reutilização de código.
 \choice Herança não afeta a reutilização de código.
 \choice Herança permite reutilizar código de uma classe filha em uma classe pai.
 \CC Herança permite reutilizar código de uma classe pai em uma classe filha.
\end{choices}

\question[1] Qual das seguintes opções é verdadeira sobre classes abstratas?
\begin{choices}
 \choice Elas podem ser instanciadas diretamente.
 \choice Elas não podem ter métodos concretos.
 \CC Elas podem ter métodos abstratos e concretos.
 \choice Elas não podem ter atributos.
\end{choices}

\question[1] Qual das seguintes opções é verdadeira sobre métodos abstratos em Java?
\begin{choices}
 \choice Eles podem ter uma implementação.
 \CC Eles não podem ter uma implementação.
 \choice Eles podem ser chamados diretamente.
 \choice Eles não podem ser declarados em interfaces.
\end{choices}

\question[1] Qual das seguintes opções é verdadeira sobre herança de classes abstratas?
\begin{choices}
 \choice Uma classe abstrata não pode herdar de outra classe abstrata.
 \choice Uma classe concreta não pode herdar de uma classe abstrata.
 \choice Uma classe abstrata não pode ter subclasses.
 \CC Uma classe abstrata pode herdar de outra classe abstrata.
\end{choices}

\question[1] Qual das seguintes opções é verdadeira sobre a implementação de métodos abstratos?
\begin{choices}
 \choice Eles devem ser implementados em uma classe abstrata.
 \choice Eles não precisam ser implementados.
 \choice Eles só podem ser implementados em interfaces.
 \CC Eles devem ser implementados em uma classe concreta que herda a classe abstrata.
\end{choices}

\question[1] Qual das seguintes opções é verdadeira sobre a criação de instâncias de classes abstratas?
\begin{choices}
 \choice Classes abstratas podem ser instanciadas diretamente.
 \CC Classes abstratas não podem ser instanciadas diretamente.
 \choice Classes abstratas só podem ser instanciadas indiretamente através de interfaces.
 \choice Classes abstratas só podem ser instanciadas se não tiverem métodos abstratos.
\end{choices}

\question[1] Qual das seguintes opções é verdadeira sobre o uso de construtores em classes abstratas?
\begin{choices}
 \CC Classes abstratas podem ter construtores.
 \choice Classes abstratas não podem ter construtores.
 \choice Construtores em classes abstratas devem ser abstratos.
 \choice Construtores em classes abstratas não podem ser chamados.
\end{choices}

\question[1] Qual das seguintes opções é verdadeira sobre a visibilidade de métodos abstratos?
\begin{choices}
 \choice Métodos abstratos devem ser privados.
 \choice Métodos abstratos devem ser protegidos.
 \CC Métodos abstratos devem ser públicos ou protegidos.
 \choice Métodos abstratos podem ter qualquer nível de visibilidade.
\end{choices}

\question[1] Qual das seguintes opções é verdadeira sobre a combinação de classes abstratas e interfaces?
\begin{choices}
 \choice Uma classe abstrata não pode implementar uma interface.
 \choice Uma interface pode herdar de uma classe abstrata.
 \choice Uma classe abstrata não pode ter métodos concretos se implementar uma interface.
 \CC Uma classe abstrata pode implementar uma interface.
\end{choices}

\question[1] Qual das seguintes opções é verdadeira sobre a implementação de interfaces em classes?
\begin{choices}
 \CC Uma classe pode implementar múltiplas interfaces.
 \choice Uma classe pode implementar apenas uma interface.
 \choice Uma interface pode implementar outra interface.
 \choice Uma interface pode herdar de uma classe.
\end{choices}

\question[1] Qual das seguintes opções é verdadeira sobre a herança de interfaces?
\begin{choices}
 \choice Uma interface não pode herdar de outra interface.
 \choice Uma interface pode herdar de uma classe.
 \CC Uma interface pode herdar de outra interface.
 \choice Uma interface não pode herdar de uma classe.
\end{choices}

\question[1] Qual das seguintes opções é verdadeira sobre a visibilidade de métodos em interfaces?
\begin{choices}
 \choice Métodos em interfaces são privados por padrão.
 \CC Métodos em interfaces são públicos por padrão.
 \choice Métodos em interfaces são protegidos por padrão.
 \choice Métodos em interfaces podem ter qualquer nível de visibilidade.
\end{choices}

\question[1] O que é polimorfismo em Java?
\begin{choices}
 \choice A capacidade de uma classe herdar de múltiplas classes.
 \choice A capacidade de um método ter múltiplas implementações.
 \CC A capacidade de uma variável referenciar objetos de diferentes tipos.
 \choice A capacidade de uma classe ter múltiplos métodos com o mesmo nome.
\end{choices}

\question[1] Qual das seguintes opções é verdadeira sobre a sobrescrita de métodos em Java?
\begin{choices}
 \choice Ela permite que uma classe filha herde múltiplos métodos com o mesmo nome de sua classe pai.
 \choice Ela permite que uma classe filha tenha múltiplos métodos com o mesmo nome, mas com diferentes assinaturas.
 \choice Ela permite que uma classe filha herde métodos estáticos de sua classe pai.
 \CC Ela permite que uma classe filha forneça uma implementação específica de um método que já está definido em sua classe pai.
\end{choices}

\question[1] Qual das seguintes opções é verdadeira sobre a sobrecarga de métodos em Java?
\begin{choices}
 \CC Ela permite que uma classe tenha múltiplos métodos com o mesmo nome, mas com diferentes assinaturas.
 \choice Ela permite que uma classe tenha múltiplos métodos com o mesmo nome e a mesma assinatura.
 \choice Ela permite que uma classe herde métodos com o mesmo nome de sua classe pai.
 \choice Ela permite que uma classe tenha métodos estáticos com o mesmo nome.
\end{choices}

\question[1] Qual das seguintes opções é verdadeira sobre o uso de polimorfismo com interfaces?
\begin{choices}
 \choice Uma variável de tipo interface pode referenciar apenas objetos de classes abstratas.
 \choice Uma variável de tipo interface não pode referenciar objetos de classes concretas.
 \CC Uma variável de tipo interface pode referenciar qualquer objeto que implemente essa interface.
 \choice Uma variável de tipo interface pode referenciar apenas objetos de uma única classe.
\end{choices}

\question[1] Qual das seguintes opções é verdadeira sobre o uso de polimorfismo com classes abstratas?
\begin{choices}
 \CC Uma variável de tipo classe abstrata pode referenciar qualquer objeto de uma classe concreta que estenda essa classe abstrata.
 \choice Uma variável de tipo classe abstrata pode referenciar apenas objetos de outras classes abstratas.
 \choice Uma variável de tipo classe abstrata não pode referenciar objetos de classes concretas.
 \choice Uma variável de tipo classe abstrata pode referenciar apenas objetos de uma única classe concreta.
\end{choices}

\question[1] Qual das seguintes opções é verdadeira sobre o uso de polimorfismo com métodos estáticos?
\begin{choices}
 \choice Métodos estáticos podem ser sobrescritos para suportar polimorfismo.
 \CC Métodos estáticos não podem ser sobrescritos e não suportam polimorfismo.
 \choice Métodos estáticos podem ser sobrecarregados para suportar polimorfismo.
 \choice Métodos estáticos podem ser herdados para suportar polimorfismo.
\end{choices}

\question[1] Qual das seguintes opções é verdadeira sobre classes em Java?
\begin{choices}
 \choice Elas são instâncias de objetos.
 \CC Elas são modelos para criar objetos.
 \choice Elas não podem ter métodos.
 \choice Elas não podem ter atributos.
\end{choices}

\question[1] Qual das seguintes opções é verdadeira sobre o uso de modificadores de acesso em classes?
\begin{choices}
 \choice Classes podem ser públicas, privadas ou protegidas.
 \choice Classes podem ser públicas ou protegidas, mas não privadas.
 \choice Classes podem ser públicas ou privadas, mas não protegidas.
 \CC Classes podem ser públicas ou ter visibilidade de pacote (default).
\end{choices}

\question[2] Qual das seguintes opções é verdadeira sobre objetos em Java?
\begin{choices}
 \choice Eles são modelos para criar classes.
 \CC Eles são instâncias de classes.
 \choice Eles não podem ter métodos.
 \choice Eles não podem ter atributos.
\end{choices}

\question[2] Qual método da classe Object é usado para comparar dois objetos para igualdade?
\begin{choices}
 \choice compareTo()
 \choice hashCode()
 \CC equals()
 \choice toString()
\end{choices}

\question[2] Qual das seguintes opções é verdadeira sobre o uso de métodos estáticos em objetos?
\begin{choices}
 \CC Métodos estáticos podem ser chamados sem criar uma instância do objeto.
 \choice Métodos estáticos não podem ser chamados sem criar uma instância do objeto.
 \choice Métodos estáticos podem acessar atributos de instância diretamente.
 \choice Métodos estáticos não podem acessar outros métodos estáticos.
\end{choices}

\question[2] Qual das seguintes opções é verdadeira sobre o uso de métodos finais em objetos?
\begin{choices}
 \choice Métodos finais podem ser sobrescritos.
 \CC Métodos finais não podem ser sobrescritos.
 \choice Métodos finais não podem ser chamados diretamente.
 \choice Métodos finais não podem ser estáticos.
\end{choices}

\question[2] Qual palavra-chave é usada para declarar um atributo ou método de classe em Java?
\begin{choices}
 \CC static
 \choice final
 \choice const
 \choice global
\end{choices}

\question[2] Qual das seguintes opções é verdadeira sobre métodos estáticos?
\begin{choices}
 \choice Eles podem acessar atributos de instância diretamente.
 \choice Eles podem ser sobrescritos.
 \CC Eles podem ser chamados sem criar uma instância da classe.
 \choice Eles não podem acessar outros métodos estáticos.
\end{choices}

\question[2] Qual das seguintes opções é verdadeira sobre atributos estáticos?
\begin{choices}
 \CC Eles são compartilhados por todas as instâncias da classe.
 \choice Eles são únicos para cada instância da classe.
 \choice Eles não podem ser modificados após a inicialização.
 \choice Eles não podem ser acessados diretamente.
\end{choices}

\question[2] Qual das seguintes opções é verdadeira sobre o uso de blocos estáticos em Java?
\begin{choices}
 \choice Blocos estáticos são executados quando uma instância da classe é criada.
 \choice Blocos estáticos não podem acessar atributos estáticos.
 \choice Blocos estáticos não podem acessar métodos estáticos.
 \CC Blocos estáticos são executados quando a classe é carregada pela primeira vez.
\end{choices}

\question[2] Qual das seguintes opções é verdadeira sobre o uso de métodos estáticos em subclasses?
\begin{choices}
 \CC Métodos estáticos não podem ser sobrescritos em subclasses.
 \choice Métodos estáticos podem ser sobrescritos em subclasses.
 \choice Métodos estáticos podem acessar atributos de instância diretamente.
 \choice Métodos estáticos não podem ser chamados diretamente.
\end{choices}

\question[2] Qual das seguintes opções é verdadeira sobre o uso de métodos estáticos privados?
\begin{choices}
 \choice Métodos estáticos privados podem ser chamados fora da classe.
 \CC Métodos estáticos privados não podem ser chamados fora da classe.
 \choice Métodos estáticos privados podem acessar atributos de instância diretamente.
 \choice Métodos estáticos privados não podem acessar outros métodos estáticos.
\end{choices}

\question[2] Qual das seguintes opções melhor descreve uma associação em Java?
\begin{choices}
 \choice Uma relação 'parte-todo' onde a parte não pode existir sem o todo.
 \choice Uma relação 'parte-todo' onde a parte pode existir independentemente do todo.
 \CC Uma relação entre duas classes onde uma classe contém uma referência a outra.
 \choice Uma relação entre duas classes onde uma classe herda da outra.
\end{choices}

\question[2] Qual das seguintes opções melhor descreve uma composição em Java?
\begin{choices}
 \CC Uma relação 'parte-todo' onde a parte não pode existir sem o todo.
 \choice Uma relação 'parte-todo' onde a parte pode existir independentemente do todo.
 \choice Uma relação entre duas classes onde uma classe contém uma referência a outra.
 \choice Uma relação entre duas classes onde uma classe herda da outra.
\end{choices}

\question[2] Qual das seguintes opções melhor descreve uma agregação em Java?
\begin{choices}
 \choice Uma relação 'parte-todo' onde a parte não pode existir sem o todo.
 \CC Uma relação 'parte-todo' onde a parte pode existir independentemente do todo.
 \choice Uma relação entre duas classes onde uma classe contém uma referência a outra.
 \choice Uma relação entre duas classes onde uma classe herda da outra.
\end{choices}

\question[2] Qual das seguintes opções é verdadeira sobre a diferença entre composição e agregação?
\begin{choices}
 \choice Na composição, a parte pode existir independentemente do todo.
 \choice Na agregação, a parte não pode existir sem o todo.
 \CC Na composição, a parte não pode existir sem o todo.
 \choice Na agregação, a parte é sempre uma instância da classe pai.
\end{choices}

\question[2] Qual das seguintes opções é verdadeira sobre a implementação de composição em Java?
\begin{choices}
 \CC A classe composta contém uma referência à classe componente e controla seu ciclo de vida.
 \choice A classe composta contém uma referência à classe componente, mas não controla seu ciclo de vida.
 \choice A classe componente contém uma referência à classe composta e controla seu ciclo de vida.
 \choice A classe componente contém uma referência à classe composta, mas não controla seu ciclo de vida.
\end{choices}

\question[2] Qual das seguintes opções é verdadeira sobre a implementação de agregação em Java?
\begin{choices}
 \choice A classe agregada contém uma referência à classe agregadora e controla seu ciclo de vida.
 \choice A classe agregada contém uma referência à classe agregadora, mas não controla seu ciclo de vida.
 \choice A classe agregadora contém uma referência à classe agregada e controla seu ciclo de vida.
 \CC A classe agregadora contém uma referência à classe agregada, mas não controla seu ciclo de vida.
\end{choices}

\question[2] O que é encapsulamento em Java?
\begin{choices}
 \choice A capacidade de uma classe herdar de múltiplas classes.
 \CC A capacidade de uma classe ocultar seus detalhes de implementação.
 \choice A capacidade de uma classe ter múltiplos métodos com o mesmo nome.
 \choice A capacidade de uma classe ter atributos e métodos estáticos.
\end{choices}

\question[2] Qual das seguintes opções é uma prática comum para implementar encapsulamento?
\begin{choices}
 \choice Tornar todos os atributos públicos.
 \choice Tornar todos os métodos privados.
 \choice Usar herança para acessar atributos privados.
 \CC Usar métodos getter e setter para acessar atributos privados.
\end{choices}

\question[2] Qual das seguintes opções é verdadeira sobre o uso de métodos getter e setter?
\begin{choices}
 \choice Eles permitem acesso direto aos atributos privados.
 \choice Eles não podem ser usados para acessar atributos privados.
 \CC Eles permitem acesso controlado aos atributos privados.
 \choice Eles são usados apenas para métodos estáticos.
\end{choices}

\question[2] Qual das seguintes opções é verdadeira sobre o uso de construtores privados?
\begin{choices}
 \CC Construtores privados impedem a criação de instâncias fora da classe.
 \choice Construtores privados permitem a criação de instâncias fora da classe.
 \choice Construtores privados não podem ser usados em classes abstratas.
 \choice Construtores privados são usados apenas em interfaces.
\end{choices}

\question[2] Qual das seguintes opções é verdadeira sobre o uso de atributos privados?
\begin{choices}
 \choice Atributos privados podem ser acessados diretamente fora da classe.
 \choice Atributos privados não podem ser usados em métodos estáticos.
 \choice Atributos privados são usados apenas em interfaces.
 \CC Atributos privados não podem ser acessados diretamente fora da classe.
\end{choices}

\question[2] Qual das seguintes opções é verdadeira sobre o uso de métodos privados?
\begin{choices}
 \choice Métodos privados podem ser chamados diretamente fora da classe.
 \CC Métodos privados não podem ser chamados diretamente fora da classe.
 \choice Métodos privados não podem ser usados em classes abstratas.
 \choice Métodos privados são usados apenas em interfaces.
\end{choices}

\question[2] Qual das seguintes opções é verdadeira sobre o uso de pacotes para encapsulamento?
\begin{choices}
 \choice Pacotes não permitem a organização de classes relacionadas.
 \choice Pacotes podem ter atributos privados.
 \CC Pacotes permitem a organização de classes relacionadas.
 \choice Pacotes podem ter métodos privados.
\end{choices}

\question[2] Qual das seguintes opções é verdadeira sobre o uso de modificadores de acesso para encapsulamento?
\begin{choices}
 \choice Modificadores de acesso não permitem controlar a visibilidade de atributos e métodos.
 \CC Modificadores de acesso permitem controlar a visibilidade de atributos e métodos.
 \choice Modificadores de acesso podem ser usados apenas em interfaces.
 \choice Modificadores de acesso podem ser usados apenas em classes abstratas.
\end{choices}

\question[2] Qual modificador de acesso permite que um atributo ou método seja acessado dentro do mesmo pacote e por subclasses?
\begin{choices}
 \choice public
 \choice private
 \choice default
 \CC protected
\end{choices}

\question[2] Qual modificador de acesso permite que um atributo ou método seja acessado apenas dentro do mesmo pacote?
\begin{choices}
 \choice public
 \choice protected
 \choice private
 \CC default
\end{choices}

\question[2] Qual das seguintes opções é verdadeira sobre o uso da palavra-chave import em Java?
\begin{choices}
 \CC Ela é usada para importar pacotes e classes.
 \choice Ela é usada para declarar pacotes.
 \choice Ela é usada para definir a visibilidade de métodos.
 \choice Ela é usada para criar herança múltipla.
\end{choices}

\question[2] Qual das seguintes opções é verdadeira sobre a estrutura de diretórios de pacotes em Java?
\begin{choices}
 \choice A estrutura de diretórios não precisa corresponder à estrutura de pacotes.
 \choice A estrutura de diretórios deve ser plana.
 \choice A estrutura de diretórios deve ser única para cada classe.
 \CC A estrutura de diretórios deve corresponder à estrutura de pacotes.
\end{choices}

\question[2] Qual das seguintes opções é verdadeira sobre o uso de pacotes padrão (default package) em Java?
\begin{choices}
 \choice Classes no pacote padrão podem ser acessadas de qualquer pacote.
 \CC Classes no pacote padrão só podem ser acessadas dentro do mesmo pacote.
 \choice Classes no pacote padrão não podem ser importadas.
 \choice Classes no pacote padrão podem ser públicas ou privadas.
\end{choices}

\question[2] Qual das seguintes opções é verdadeira sobre o uso de pacotes para modularização?
\begin{choices}
 \CC Pacotes permitem a modularização do código.
 \choice Pacotes impedem a modularização do código.
 \choice Pacotes não afetam a modularização do código.
 \choice Pacotes são usados apenas para definir a visibilidade de métodos.
\end{choices}

\question[2] Qual das seguintes opções é verdadeira sobre o uso de pacotes para evitar conflitos de nomes?
\begin{choices}
 \choice Pacotes não ajudam a evitar conflitos de nomes entre classes.
 \choice Pacotes são usados apenas para definir a visibilidade de métodos.
 \CC Pacotes ajudam a evitar conflitos de nomes entre classes.
 \choice Pacotes são usados apenas para agrupar classes relacionadas.
\end{choices}


\end{multicols}        
     \end{questions}	
\printmyanswers
\end{document}

